\section{Funzione di valutazione}

Uno degli aspetti chiave di un motore scacchistico è
avere una funzione che valuti staticamente una posizione senza esplorarla.

\subsection{Definizione matematica}

$$ \text{evaluate}: \text{Board} \longrightarrow \mathbb{Z} $$

\subsubsection{Input:}
L'input è un oggetto \texttt{Board} contenente:
\begin{itemize}
\item Le otto righe della scacchiera (disposizione pezzi)
\item Informazioni supplementari: en passant, diritti di arrocco, turno, contatori mosse
\end{itemize}

\subsubsection{Output:}
Un numero intero (centipawns) che rappresenta la valutazione della posizione:
\begin{itemize}
\item \textbf{Valore positivo}: Vantaggio per il bianco
\item \textbf{Valore negativo}: Vantaggio per il nero
\item \textbf{Zero}: Posizione equilibrata
\item \textbf{$\pm\infty$}: Scacco matto (vittoria/sconfitta)
\end{itemize}

Scala tipica: 100 centipawns = valore di un pedone.

\subsection{Implementazione della funzione}

La funzione \texttt{Engine::evaluate(const chess::Board\& board)} implementa
la valutazione attraverso diversi componenti:

\subsubsection{1. Riconoscimento scacco matto}

\begin{lstlisting}[language=C++]
if (board.isCheckmate(board.getActiveColor())) {
    return evaluateCheckmate(board);
}
\end{lstlisting}

Se la posizione è scacco matto, restituisce immediatamente:
\begin{itemize}
\item \texttt{POS\_INF} se il bianco ha vinto (nero è matto)
\item \texttt{NEG\_INF} se il nero ha vinto (bianco è matto)
\end{itemize}

\subsubsection{2. Delta materiale}

Calcolo del vantaggio materiale usando la funzione \texttt{getMaterialDeltaFAST()}:

\begin{lstlisting}[language=C++]
int64_t materialDelta = getMaterialDeltaFAST(board);
eval += materialDelta;
\end{lstlisting}

\textbf{Valori dei pezzi} (definiti in \texttt{basicrules.hpp}):
\begin{center}
\begin{tabular}{|c|c|}
\hline
\textbf{Pezzo} & \textbf{Valore (centipawns)} \\
\hline
Pedone   & 100 \\
Cavallo  & 320 \\
Alfiere  & 330 \\
Torre    & 500 \\
Regina   & 900 \\
Re       & 20000 (valore nominale) \\
\hline
\end{tabular}
\end{center}

\textbf{Implementazione efficiente}:
\begin{lstlisting}[language=C++]
int w = __builtin_popcountll(b.pawns_bb[0]);
int bl = __builtin_popcountll(b.pawns_bb[1]);
delta += (w - bl) * PIECE_VALUES[PAWN];
\end{lstlisting}

Usa \texttt{popcount} hardware per contare pezzi nei bitboard in tempo O(1).

\subsubsection{3. Riconoscimento fase di gioco}

La valutazione distingue tra \textbf{mediogioco} ed \textbf{endgame}:

\begin{lstlisting}[language=C++]
int nonPawnNonKingPieces = __builtin_popcountll(minorMajors);
bool isEndgame = (nonPawnNonKingPieces <= PHASE_FINAL_THRESHOLD);
\end{lstlisting}

Conta pezzi pesanti (cavalli, alfieri, torri, regine):
\begin{itemize}
\item Se $\leq$ \texttt{PHASE\_FINAL\_THRESHOLD}: è endgame
\item Altrimenti: è mediogioco/apertura
\end{itemize}

Questo influenza:
\begin{itemize}
\item Piece-Square Tables del re (centralizzazione in endgame)
\item Piece-Square Tables dei pedoni (valore diverso in endgame)
\end{itemize}

\subsubsection{4. Piece-Square Tables (PSQT)}

Bonus/penalty in base alla posizione di ogni pezzo sulla scacchiera.

\textbf{Tavole disponibili} (in \texttt{piecevaluetables.hpp}):
\begin{itemize}
\item \texttt{PAWN\_VALUES\_TABLE}: Pedoni nel mediogioco
\item \texttt{PAWN\_END\_GAME\_VALUES\_TABLE}: Pedoni nell'endgame
\item \texttt{KNIGHT\_VALUES\_TABLE}: Cavalli (preferiscono il centro)
\item \texttt{BISHOP\_VALUES\_TABLE}: Alfieri (preferiscono diagonali lunghe)
\item \texttt{ROOK\_VALUES\_TABLE}: Torri (preferiscono colonne aperte)
\item \texttt{QUEEN\_VALUES\_TABLE}: Regina (centro/mobilità)
\item \texttt{KING\_MIDDLE\_GAME\_VALUES\_TABLE}: Re nel mediogioco (preferisce angoli)
\item \texttt{KING\_END\_GAME\_VALUES\_TABLE}: Re nell'endgame (preferisce centro)
\end{itemize}

\textbf{Algoritmo}:
\begin{lstlisting}[language=C++]
auto addPsqtFor = [&](uint64_t bbWhite, uint64_t bbBlack,
                       auto valueSelectorWhite) {
    // Iterazione sui pezzi bianchi
    uint64_t bb = bbWhite;
    while (bb) {
        uint8_t sq = __builtin_ctzll(bb);  // trova bit meno significativo
        bb &= (bb - 1);                    // rimuovi quel bit
        eval += valueSelectorWhite(sq);
    }

    // Iterazione sui pezzi neri (con indice specchiato)
    bb = bbBlack;
    while (bb) {
        uint8_t sq = __builtin_ctzll(bb);
        bb &= (bb - 1);
        uint8_t idx = mirrorIndex(sq);
        eval -= valueSelectorWhite(idx);
    }
};
\end{lstlisting}

\textbf{Mirror index}: I neri sono valutati come se fossero dalla loro prospettiva.
La funzione \texttt{mirrorIndex(sq)} riflette la casa verticalmente: $\text{mirror}(sq) = 56 - (sq \% 8) + (sq / 8)$.

\subsubsection{5. Bonus arrocco}

\begin{lstlisting}[language=C++]
if (castled) eval += isWhite ? CASTLING_BONUS : -CASTLING_BONUS;
\end{lstlisting}

Rileva configurazioni di pezzi che indicano arrocco completato:
\begin{itemize}
\item \textbf{Bianco O-O}: Re su g1, torre su f1
\item \textbf{Bianco O-O-O}: Re su c1, torre su d1
\item \textbf{Nero O-O}: Re su g8, torre su f8
\item \textbf{Nero O-O-O}: Re su c8, torre su d8
\end{itemize}

Valore: \texttt{CASTLING\_BONUS} (definito in \texttt{basebonuspenaltyvalues.hpp}).

\subsection{Criteri di valutazione implementati}

\begin{enumerate}
\item \textbf{Materiale}: Differenza in valore dei pezzi
\item \textbf{Posizione pezzi}: Piece-Square Tables (centralizzazione, controllo)
\item \textbf{Fase di gioco}: Adattamento apertura/mediogioco vs endgame
\item \textbf{Sicurezza del re}: Re protetto (angolo) nel mediogioco, attivo (centro) nell'endgame
\item \textbf{Arrocco}: Bonus per arrocco completato
\item \textbf{Struttura pedonale}: Pedoni avanzati valgono di più (via PSQT)
\end{enumerate}

\subsection{Criteri non ancora implementati (TODO)}

Dalla sezione \texttt{todo.tex}, fattori da considerare in future iterazioni:

\subsubsection{Generale:}
\begin{itemize}
\item Mobilità pezzi (numero mosse disponibili)
\item King's safety (pedoni davanti al re)
\item Pawn islands (gruppi di pedoni separati)
\item Pedoni doppi/tripli (penalità)
\item Pedoni passati (bonus)
\item Pedoni in prossimità di promozione (7° e 2° traversa)
\item Pezzi inchiodati (pinned pieces)
\item Occupazione delle ali
\item Alfieri in fianchetto
\item Coppia alfieri (bishop pair bonus)
\item Torri su colonne aperte/semi-aperte
\end{itemize}

\subsubsection{Endgame specifico:}
\begin{itemize}
\item Pedoni passati supportati da altri pedoni
\item Opposizione del re
\item Zugzwang
\item Regola del quadrato per pedoni
\end{itemize}

\subsection{Performance della valutazione}

La funzione \texttt{evaluate()} è chiamata milioni di volte durante la ricerca, quindi
l'efficienza è critica.

\textbf{Ottimizzazioni implementate}:
\begin{itemize}
\item \textbf{Bitboard operations}: Uso di popcount e ctz hardware
\item \textbf{Lambda inline}: Il compilatore può espandere le lambda inline
\item \textbf{Lookup tables}: PSQT pre-calcolate, nessun calcolo a runtime
\item \textbf{Early exit}: Scacco matto rilevato immediatamente
\item \textbf{Cache-friendly}: Bitboard contigui in memoria
\end{itemize}

\textbf{Tempo tipico}: ~50-100 ns per valutazione su hardware moderno.

\subsection{Funzione alternativa: getMaterialDelta()}

Esiste una versione alternativa del calcolo materiale:

\begin{lstlisting}[language=C++]
int64_t Engine::getMaterialDelta(const chess::Board& b) noexcept {
    // Usa un polinomio per calcolare valori pezzi
    auto pieceValue = [](uint8_t x) {
        return (x * (-134220 + x * (304540 + ...))) / 36.0;
    };
    // ...
}
\end{lstlisting}

\textbf{Differenza}:
\begin{itemize}
\item Itera su tutte le 64 case invece di usare bitboard
\item Usa un polinomio interpolante per calcolare valori
\item Più lenta di \texttt{getMaterialDeltaFAST()}
\item Mantenuta per compatibilità o testing
\end{itemize}

\textbf{Versione usata attualmente}: \texttt{getMaterialDeltaFAST()} per performance ottimali.


